\documentclass{report}
\usepackage{amsmath,amssymb}
\setlength{\parindent}{0mm}
\setlength{\parskip}{1em}
\begin{document}
\begin{center}
\rule{6in}{1pt} \
{\large Erik G. Boman \\
{\bf ShyLU: A new hybrid direct-iterative preconditioner and solver for sparse linear systems.}}

Sandia National Labs \\ PO Box 5800 \\ Albuquerque \\ NM 87185-1318
\\
{\tt egboman@sandia.gov}\\
Sivasankaran Rajamanickam\\
Michael A. Heroux\\
Heidi Thornquist\end{center}

We present ShyLU, a new hybrid direct-iterative preconditioner and solver
for general sparse linear systems.
We use a Schur complement approach where subproblems are solved using a
direct solver while the Schur complement is solved iteratively.
We employ a block decomposition, which is formed in a fully automatic way
using state-of-the-art partitioning and ordering algorithms in
the Zoltan and Scotch software packages. We discuss the choice of
partitioning strategy, and in particular, narrow vs. wide separators,
and study the impact on the Schur complement system.

ShyLU was developed in Trilinos and will become a Trilinos package.
We have two main goals: First, we wish ShyLU to be a fast and robust
solver/preconditioner for multicore compute nodes. Second, due to its
modular design, ShyLU can serve as an research platform to compare
algorithms. For example, ShyLU supports two different approaches
to forming a sparse preconditioner for the Schur complement:
dropping and probing.

We show results that demonstrate the robustness and effectiveness
of ShyLU for problems from a variety of applications. We have found
it is particularly effective on circuit simulation problems where
many iterative methods fail.

We also show results from a hybrid MPI + threads version of ShyLU.
We observe that on multicore systems with many cores per node, it
is advantageous to use MPI even within a single compute node.

Finally, we describe how ShyLU in future work can be used as a
subdomain solver within a domain decomposition framework. This
will help overcome scalability issues in our Schur complement approach
for very large problems.


\end{document}
